% 霍普夫代数（综述）
% license CCBYSA3
% type Wiki

本文根据 CC-BY-SA 协议转载翻译自维基百科\href{https://en.wikipedia.org/wiki/Hopf_algebra}{相关文章}

在数学中，霍普夫代数（Hopf algebra，以海因茨·霍普夫 Heinz Hopf 命名）是一种结构，它既是一个（含幺结合的）代数，又是一个（含余元的余结合）余代数；这些结构的相容性使它成为一个双代数，此外它还配备了一个满足特定性质的反同态。霍普夫代数的表示论尤其优美，因为兼容的余乘法、余元以及反元的存在，使得能够构造表示的张量积、平凡表示以及对偶表示。

霍普夫代数自然地出现在代数拓扑中（其起源之处，并与 H-空间概念相关）、群概形理论、群论（通过群环的概念），以及许多其他地方，因此可能是最常见的一类双代数。霍普夫代数也作为独立的研究对象受到关注，一方面有许多关于具体类实例的研究，另一方面也存在分类问题。它们拥有广泛的应用，从凝聚态物理与量子场论，到弦论以及大型强子对撞机（LHC）现象学。
\subsection{形式定义}
形式上，霍普夫代数是域$K$上的一个（结合且余结合的）双代数$H$，并配备一个$K$-线性映射$S: H \to H$（称为反元，antipode），使得下列图表交换：
